\documentclass[11pt]{article}
\usepackage[a4paper,margin=18mm]{geometry}
\usepackage{fontspec}
\setmainfont{Latin Modern Roman}
\usepackage{amsmath,amssymb,amsthm,amscd,graphicx,color}
\usepackage{xurl}
\usepackage[unicode,hidelinks]{hyperref}
\allowdisplaybreaks
\setlength\emergencystretch{3em}

\makeatletter
\newtheorem{theorem}{Theorem}
\makeatother
\usepackage{amscd}
\usepackage[all]{xy}

\numberwithin{equation}{section}

\newtheorem{Theorem}{Theorem}[section]
\newtheorem{Corollary}[Theorem]{Corollary}
\newtheorem{Lemma}[Theorem]{Lemma}
\newtheorem{Proposition}[Theorem]{Proposition}
\newtheorem{Conjecture}[Theorem]{Conjecture}
 { \theoremstyle{definition}
\newtheorem{Definition}[Theorem]{Definition}
\newtheorem{Note}[Theorem]{Note}
\newtheorem{Example}[Theorem]{Example}
\newtheorem{Remark}[Theorem]{Remark}
\newtheorem{Condition}[theorem]{Condition}
\newtheorem{Notation}[theorem]{Notation}
\newtheorem{Exercise}[theorem]{Exercise}
\newtheorem{Question}[theorem]{Question} }

\newcommand{\NN}{{\mathbb{N}}} %natural numbers
\newcommand{\QQ}{{\mathbb{Q}}} %rational numbers
\newcommand{\RR}{{\mathbb{R}}} %real numbers or cartesian space
\newcommand{\TT}{{\mathbb{T}}} %torus
\newcommand{\ZZ}{{\mathbb{Z}}} %integers
\newcommand{\Hom}{{\operatorname{Hom}}} %Hom
\newcommand{\inter}{{\operatorname{int}}} %interior
\newcommand{\pr}{{\operatorname{pr}}} %projection
\newcommand{\supp}{{\operatorname{supp}}} %support
\newcommand{\g}{{\mathfrak{g}}} %generic Lie algebra
\newcommand{\GL}{{\operatorname{GL}}} %general linear
\newcommand{\lie}[1]{{\mathfrak{\lowercase{#1}}}} %Lie algebra of a Lie group
\renewcommand{\O}{{\operatorname{O}}} %orthogonal
\newcommand{\SL}{{\operatorname{SL}}} %special linear
\newcommand{\SO}{{\operatorname{SO}}} %special orthogonal
\newcommand{\Sp}{{\operatorname{Sp}}} %symplectic linear
\newcommand{\SU}{{\operatorname{SU}}} %special unitary
\newcommand{\U}{{\operatorname{U}}} %unitary

\newcommand{\CIN}{{C^\infty}} %infinitely differentiable
\newcommand{\hook}{{\lrcorner\,}} %hook or interior multiplication
\def \red {{\text{red}}}
\def \reg {{\text{reg}}}
\def \ol {\overline}
\newcommand{\toto}{{\rightrightarrows}}
\def \calD {{\mathcal{D}}}
\def \calO {{\mathcal{O}}}
\def \tpsi {{\widetilde{\psi}}}
\def \tU {{\widetilde{U}}}
\DeclareMathOperator {\image}{Image}

\newcommand{\hypref}[2]{{#2~\ref{#1}}}



\begin{document}
\begin{center}\Large Diffeological differential forms on orbit spaces equal basic forms when the Lie group identity component acts properly\end{center}
\noindent\textbf{Stable ID:} 1893\par
\noindent\textbf{Authors:} Yael Karshon; Jordan Watts\par
\noindent\textbf{Original work:} Basic Forms and Orbit Spaces: a Diffeological Approach\par
\noindent\textbf{Version:} Official SIGMA 12 (2016) article 026, DOI 10.3842/SIGMA.2016.026; journal PDF and native TeX acquired 2026-10-01, exact bytes pinned by SHA256\par
\noindent\textbf{Selected task scope:} Smooth Lie group G action on manifold M, quotient M/G with quotient diffeology and form spaces with standard functional diffeology. General part(i): pullback injects into G-basic forms and is an isomorphism of differential graded algebras and diffeological diffeomorphism onto its image. Part(ii): if G0 acts properly, the image is all invariant horizontal basic forms. No condition that G itself act properly/free or that M/G be Hausdorff.\par
\noindent\textbf{Difficulty:} H2: The quotient may have singular stabilizers and be non-Hausdorff, so pointwise equality of orbits does not provide a smooth group-valued relation between plots. The proof combines a countable-component Baire argument with slice induction on group dimension and a fixed-point velocity comparison for compact connected stabilizers. It also proves the pullback inverse is smooth in the functional diffeology, beyond the classical free proper quotient argument.\par
\noindent\textbf{Editorial boundary note (not author text):} The complete original Theorem 1.1 and full author Sections 3–5 proof are retained, with all quotient/functional diffeology and form definitions. New pullback injection/basicity, smooth functional-diffeology inverse, quotient-in-stages correspondence, component-coset closedness/Baire density, fixed-point averaged velocity comparison, slice associated-bundle reduction and group-dimension induction are supplied as one main descent outcome. The plot criterion and proper slice theorem remain cited established prerequisites. Properness is required only for the identity component; no freeness, full-group properness or Hausdorff quotient is added. Orbifold and symplectic-quotient applications are unselected.\par
\noindent\textbf{Source:} \url{https://sigma-journal.com/2016/026/}\quad DOI: \url{https://doi.org/10.3842/SIGMA.2016.026}\par
\noindent\textbf{License and attribution:} Copyright retained by the named authors. Published by SIGMA under CC BY 4.0: \url{https://creativecommons.org/licenses/by/4.0/}. Source: \url{https://sigma-journal.com/about.html}. Adaptation consists of standalone excerpt formatting, cover and numbering restoration; original author language and mathematical text are retained.\par
\medskip\hrule\medskip\noindent\textbf{Author text begins below.}\par
\setcounter{section}{0}
\setcounter{subsection}{0}
\setcounter{subsubsection}{0}
\setcounter{Theorem}{0}
\setcounter{equation}{0}
% Original source lines 109--153
\section{Introduction}\label{sec:intro}

Let $M$ be a smooth manifold and $G$ a Lie group acting on $M$.
A \emph{basic differential form} on $M$
is a dif\/ferential form that is $G$-invariant and horizontal;
the latter means that
evaluating the form on any vector that is tangent to a $G$-orbit yields $0$.
Basic dif\/ferential forms constitute a~subcomplex of the de Rham complex.
If $G$ acts properly and with a constant orbit-type, then the quotient~$M/G$ is a manifold,
and, denoting the quotient map by $\pi \colon M \to M/G$,
the pullback by this map
gives an isomorphism of the de Rham complex on~$M/G$
with the complex of basic forms on~$M$.
Even if~$M/G$ is not a manifold, if $G$ acts properly, then
the cohomology of the complex of basic forms
is isomorphic to the singular cohomology of $M/G$ with real coef\/f\/icients;
this was shown by Koszul in 1953~\cite{koszul} for compact group actions
and by Palais in 1961~\cite{palais} for proper group actions.
In light of these facts, some authors def\/ine the de Rham complex on~$M/G$
to be the complex of basic forms on~$M$.

There is another, intrinsic, def\/inition of a dif\/ferential form on $M/G$,
which comes from viewing $M/G$ as a dif\/feological space (see Section~\ref{sec:diffeology}). This def\/inition agrees with the usual one when $M/G$ is a manifold. Dif\/ferential forms on dif\/feological spaces
admit exterior derivatives,
wedge products, and pullbacks under smooth maps.
Spaces of dif\/ferential forms are themselves dif\/feological spaces too.
With this notion, here is our main result:

\begin{Theorem}\label{t:main}\looseness=-1
Let $G$ be a Lie group acting on a manifold $M$.
Let $\pi\colon M\to M/G$ be the quotient map.
\begin{enumerate}\itemsep=0pt
\item[$(i)$] The pullback map $\pi^* \colon \Omega^*(M/G) \to \Omega^*(M)$
is one-to-one. Its image is contained in the space $\Omega^*_{\rm basic}(M)$
of basic forms. As a map to its image,
the map $\pi^*$ is an isomorphism of differential graded algebras
and a diffeological diffeomorphism.
\item[$(ii)$] If the restriction of the action to the identity component of~$G$
is proper, then the image of the pullback map is equal to the space
of basic forms:
\begin{gather*}\xymatrix{
\pi^* \colon \ \Omega^*(M/G) \ar[r]^(.55){\cong}
& \Omega^*_{\rm basic}(M)}.\end{gather*}
\end{enumerate}
\end{Theorem}

\setcounter{section}{1}
\setcounter{subsection}{0}
\setcounter{subsubsection}{0}
\setcounter{Theorem}{2}
\setcounter{equation}{0}
% Original source lines 279--1398
\section{Background on dif\/feological spaces}\label{sec:diffeology}

In this section we review the basics of dif\/feology, dif\/feological dif\/ferential forms, and Lie group actions in the context of dif\/feology. For more details (e.g., the quotient dif\/feology is in fact a~dif\/feology), see Iglesias-Zemmour~\cite{iglesias}.

\subsection*{The basics of dif\/feology}

This subsection contains a review of the basics of dif\/feology; in particular, the def\/inition of a~dif\/feology and dif\/feologically smooth maps, as well as various constructions in the dif\/feological category.

\begin{Definition}[dif\/feology]
Let $X$ be a set. A \emph{parametrisation} on $X$ is a
function $p \colon U \to X$ where $U$ is an open subset of $\RR^n$
for some $n$.
A \emph{diffeology} $\mathcal{D}$ on $X$ is a set of parametrisations
that satisf\/ies the following three conditions.

\begin{enumerate}\itemsep=0pt
\item (\emph{Covering})
For every point $x \in X$ and every non-negative integer $n \in \NN$,
the constant function $p \colon \RR^n \to \{ x \} \subseteq X$
is in $\mathcal{D}$.
\item (\emph{Locality})
Let $p \colon U \to X$ be a parametrisation such that
for every point in $U$ there exists an open neighbourhood $V$ in $U$
such that $p|_V \in \mathcal{D}$. Then $p\in\mathcal{D}$.
\item (\emph{Smooth compatibility})
Let $(p \colon U\to X) \in \mathcal{D}$.
Then for every $n \in \NN$, every open subset $V \subseteq \RR^n$,
and every smooth map $F \colon V \to U$, we have $p \circ F \in \mathcal{D}$.
\end{enumerate}

A set $X$ equipped with a dif\/feology $\mathcal{D}$ is called
a \emph{diffeological space} and is denoted by $(X,\mathcal{D})$.
When the dif\/feology is understood, we drop the symbol $\mathcal{D}$.
The elements of $\calD$ are called \emph{plots}.
\end{Definition}

\begin{Example}[standard dif\/feology on a manifold]
Let $M$ be a manifold. The \emph{standard diffeology} on $M$ is the
set of all smooth maps to $M$ from open subsets of $\RR^n$
for all $n \in \NN$.
\end{Example}

\begin{Definition}[dif\/feologically smooth maps]
Let $X$ and $Y$ be two dif\/feological spaces,
and let $F \colon X\to Y$ be a map. We say that $F$ is
\emph{(diffeologically) smooth} if for any plot $p \colon U \to X$ of $X$
the composition $ F \circ p \colon U \to Y$ is a plot of $Y$.
Denote by $\CIN(X,Y)$ the set of all smooth maps
from $X$ to $Y$.
Denote by $\CIN(X)$ the set of all smooth maps from $X$ to~$\RR$,
where $\RR$ is equipped with its standard dif\/feology.
\end{Definition}

\begin{Remark}[plots]
A parametrisation is dif\/feologically smooth if and only if it is a plot.
\end{Remark}

\begin{Remark}[smooth maps between manifolds]
A map between two manifolds is dif\/feologically smooth
if and only if it is smooth in the usual sense.
In particular, if $M$ is a manifold then $\CIN(M)$ is the usual set
of smooth real valued functions.
\end{Remark}

\begin{Remark}
Dif\/feological spaces, along with dif\/feologically smooth maps,
form a category. It is shown in \cite[Theorem~3.2]{BH} that this category
is a complete and cocomplete quasi-topos.
In particular, it is closed under passing to arbitrary quotients,
subsets, function spaces, products, and coproducts.
\end{Remark}

\begin{Definition}[quotient dif\/feology]
Let $X$ be a dif\/feological space,
and let $\sim$ be an equivalence relation on $X$.
Let $Y = X/\!\!\sim$ be the quotient set,
and let $\pi \colon X\to Y$ be the quotient map.
We def\/ine the \emph{quotient diffeology} on $Y$
to be the dif\/feology for which the plots are those maps $p \colon U\to Y$
such that for every point in $U$
there exist an open neighbourhood $V \subseteq U$
and a plot $q \colon V\to X$ such that $p|_V = \pi \circ q$.
\end{Definition}

\begin{Remark}[quotient map]
Let $X$ be a dif\/feological space and $\sim$ an equivalence relation on~$X$.
Then the quotient map $\pi \colon X \to X / \!\!\sim$ is smooth.
\end{Remark}

A special case that is important to us is the quotient
of a manifold by the action of a Lie group.

\begin{Definition}[subset dif\/feology]
Let $X$ be a dif\/feological space, and let~$Y$ be a subset
of~$X$. The \emph{subset diffeology} on~$Y$ consists
of those maps to~$Y$ whose composition with the inclusion map $Y \to X$
are plots of~$X$.
\end{Definition}

\begin{Definition}[product dif\/feology]
Let $X$ and $Y$ be two dif\/feological spaces.
The \emph{product diffeology} on the set $X \times Y$ is def\/ined as follows.
Let $\operatorname{pr}_X \colon X \times Y \to X$
and $\operatorname{pr}_Y \colon X \times Y \to Y$ be the natural projections.
A parametrisation $p \colon U \to X \times Y$ is a plot
if $\operatorname{pr}_X \circ p$ and $\operatorname{pr}_Y \circ p$
are plots of $X$ and $Y$, respectively.
\end{Definition}

\begin{Definition}[standard functional dif\/feology on maps]
\label{d:functional}
Let $Y$ and $Z$ be dif\/feological spaces.
The \emph{standard functional diffeology} on $\CIN(Y,Z)$ is def\/ined as follows.
A parametrisation $p \colon U \to \CIN(Y,Z)$ is a plot if the map
 \begin{gather*} U \times Y\to Z \qquad \text{given by} \quad (u,y)\mapsto p(u)(y) \end{gather*}
is smooth.
\end{Definition}

\subsection*{Dif\/feological dif\/ferential forms}

In this subsection we review dif\/ferential forms on dif\/feological spaces, as well as introduce Proposition~\ref{p:pullbackimgchar} which is a~simple criterion crucial to the proof of Theorem~\ref{t:main}.

\begin{Definition}[dif\/ferential forms]
\label{d:diff forms}
Let $(X,\mathcal{D})$ be a dif\/feological space.
A \emph{$($diffeological$)$ differential $k$-form} $\alpha$ on $X$
is an assignment to each plot $( p \colon U \to X) \in \mathcal{D}$
a dif\/ferential $k$-form $\alpha(p) \in \Omega^k(U)$
satisfying the following \emph{smooth compatibility} condition:
for every open subset $V$ of a~Euclidean space
and every smooth map $F \colon V \to U$,
\begin{gather*} \alpha(p\circ F) = F^* (\alpha(p)) .\end{gather*}
Denote the set of dif\/ferential $k$-forms on $X$ by $\Omega^k(X)$.
\end{Definition}

\begin{Definition}[pullback map]
\label{d:pullback}
Let $X$ and $Y$ be dif\/feological spaces, and let $F \colon X\to Y$
be a dif\/feologically smooth map.
Let $\alpha$ be a dif\/ferential $k$-form on $Y$.
Def\/ine the \emph{pullback} $F^*\alpha$ to be the $k$-form on $X$
that satisf\/ies the following condition:
for every plot $p \colon U \to X$,
\begin{gather*} F^* \alpha(p) = \alpha(F\circ p) .\end{gather*}
\end{Definition}

\begin{Example} \label{e:forms-on-mflds}
Let $\alpha$ be a dif\/ferential form on a manifold $M$.
Then $(p \colon U \to M) \mapsto p^* \alpha$
def\/ines a dif\/feological dif\/ferential form on $M$.
In this way, we get an identif\/ication of the ordinary dif\/ferential forms
on $M$ with the dif\/feological dif\/ferential forms on $M$.

Let $X$ be a dif\/feological space, $\alpha$ a dif\/ferential form on $X$,
and $p \colon U \to X$ a plot.
The above identif\/ication of ordinary dif\/ferential forms on $U$
with dif\/feological dif\/ferential forms on $U$
gives $\alpha(p) = p^* \alpha$.
Henceforth, we may write $p^*\alpha$ instead of $\alpha(p)$.

\end{Example}

\begin{Example} \label{e:Omega0}
The space $\Omega^0(X)$ of dif\/feological $0$-forms
is identif\/ied with the space $\CIN(X)$
of smooth real valued functions,
by identifying the function $f$
with the $0$-form $(p \colon U \to X) \mapsto f \circ p$.
See \cite[Section~6.31]{iglesias}.
With this identif\/ication, the pullback of $0$-forms by a smooth map~$F$
becomes the precomposition of smooth real-valued functions by~$F$.
\end{Example}

\begin{Remark}[pullback is linear]\label{r:linear}
The space of dif\/ferential forms on a dif\/feological space is naturally
a linear vector space:
for $\alpha,\beta \in \Omega^k(X)$ and $a,b \in \RR$,
we def\/ine $a\alpha+b\beta \colon (p \colon U \to X) \mapsto
a \alpha(p) + b \beta(p)$.
If $F \colon X \to Y$ is a smooth map of dif\/feological spaces,
then the pullback map $F^*	 \colon \Omega^k(Y) \to \Omega^k(X)$ is linear.
\end{Remark}

\begin{Remark}[wedge product and exterior derivative]
\label{alg iso}
Let $X$ be a dif\/feological space.
Def\/ine the \emph{wedge product}
of $ \alpha \in \Omega^k(X) $ and $ \beta \in \Omega^l(X) $
to be the $(k+l)$-form
$\alpha \wedge \beta \colon (p \colon U \to X)
 \mapsto p^* \alpha \wedge p^* \beta$.
Def\/ine the \emph{exterior derivative} of $\alpha$
to be the $k+1$ form
$d\alpha \colon (p \colon U \to X) \mapsto d (p^*\alpha)$.
Then $\Omega^*(X)=\bigoplus_{k=0}^\infty\Omega^k(X)$
is a dif\/ferential graded algebra.
In particular, $\Omega^*(X)$ is an exterior algebra,
and $(\Omega^*(X),d)$ is a complex.
If $F \colon X \to Y$ is a smooth map of dif\/feological spaces,
then the pullback map $F^* \colon \Omega^*(Y) \to \Omega^*(X)$
is a morphism of dif\/ferential graded algebras;
in particular,
it intertwines the wedge products and the exterior derivatives.
\end{Remark}

\begin{Definition}[standard functional dif\/feology on forms]
\label{d:D on forms}
Let $(X,\calD)$ be a dif\/feological space.
The \emph{standard functional diffeology} on $\Omega^k(X)$
is def\/ined as follows.
A parametrisation $p \colon U \to \Omega^k(X)$ is a plot
if for every plot $(q \colon V \to X) \in \mathcal{D}$,
where $V$ is open in $\RR^n$, the map $U \times V \to \bigwedge^k \RR^n$
sending $(u,v)$ to $q^*(p(u))|_v$ is smooth.
See \cite[Section~6.29]{iglesias} for a proof that this is indeed a~dif\/feology.
\end{Definition}

\begin{Remark}\label{r:formfacts}
Let $X$ be a dif\/feological space.
We have the following facts.
\begin{enumerate}\itemsep=0pt

\item %\label{f:Omega0}
Under the identif\/ication of \hypref{e:Omega0}{Example}
of the space of dif\/feological $0$-forms
with the space of smooth real valued functions,
Def\/initions~\ref{d:functional} and~\ref{d:D on forms}
of the standard functional dif\/feology on these spaces agree.

\item \label{f:pullback is sm}
If $F \colon X \to Y$ is a smooth map to another dif\/feological space,
then the pullback map $F^* \colon \Omega^k(Y) \to \Omega^k(X)$
is smooth with respect to the standard functional dif\/feologies
on the sets of dif\/ferential forms.
See \cite[Section~6.32]{iglesias}.

\item
The exterior derivative
$d \colon \Omega^k(X)\to\Omega^{k+1}(X)$ and the wedge product
$ \Omega^k (X) \times \Omega^l (X) \to \Omega^{k+l} (X)$ are smooth.
See Sections~6.34 and~6.35 of~\cite{iglesias}.
\end{enumerate}
\end{Remark}

\begin{Proposition}[pullbacks of quotient dif\/feological forms]\label{p:pullbackimgchar}
Let $G$ be a Lie group, acting on a manifold $M$,
and let $\pi \colon M \to M/G$ be the quotient map.
Then a differential form $\alpha$ on $M$ is in the image of $\pi^*$
if and only if, for every two plots $p_1 \colon U\to M$
and $p_2 \colon U \to M$ such that $\pi \circ p_1 = \pi \circ p_2$, we have
\begin{gather*} p_1^* \alpha = p_2^* \alpha. \end{gather*}
\end{Proposition}

\begin{proof}
This result is a special case of \cite[Section~6.38]{iglesias}.
\end{proof}

\subsection*{Group actions}

Here we highlight some facts about Lie group actions in connection to dif\/feology, as well as review the def\/inition of basic forms. Finally, we introduce another crucial ingredient to the proof of Theorem~\ref{t:main}, the slice theorem (Theorem~\ref{t:slice theorem}). Koszul \cite{koszul} proved the slice theorem for compact Lie group actions, and
Palais \cite{palais} proved it for proper Lie group actions.
The proof is also described in Theorem~2.3.3 of~\cite{DK}
and in Appendix~B of~\cite{GGK}.

\begin{Lemma} \label{facts}
Let a Lie group $G$ act on a manifold $M$.
Let $x$ be a point in $M$ and $H$ its stabiliser. Then
\begin{itemize}\itemsep=0pt
\item
There exists a unique manifold structure on the quotient $G/H$
such that the quotient map $G \to G/H$ is a submersion.
The standard diffeology on this manifold
agrees with the quotient diffeology induced from $G$.

\item
There exists a unique manifold structure on the orbit $G \cdot x$
such that the inclusion map $G \cdot x \to M$ is an immersion.
The standard diffeology on this manifold agrees with the subset diffeology
induced from~$M$.

\item
The orbit map $a \mapsto a \cdot x$ from $G$ to $M$
descends to a diffeomorphism from $G/H$ to $G \cdot x$.

\item
The tangent space $T_x(G \cdot x)$
is the space of vectors $\xi_M|_x$ for $\xi \in \g$,
where $\xi_M$ is the vector field on $M$ that is induced
by the Lie algebra element $\xi$.
This space is also the image of the differential at the identity
of the orbit map $a \mapsto a \cdot x$ from $G$ to $M$.
\end{itemize}
\end{Lemma}

\begin{proof}
See \cite[Section~2, Paragraph~1]{IK}.
\end{proof}

\begin{Definition}[basic forms] \label{d:basic form}
Let $G$ be a Lie group acting on a manifold $M$.
A dif\/ferential form $\alpha$ on $M$ is \emph{horizontal}
if for any $x \in M$ and $v \in T_x(G\cdot x)$ we have
\begin{gather*} v \hook \alpha = 0.\end{gather*}
(Recall that $v\hook\alpha=\alpha(v,\cdot,\dots,\cdot)$.)
A form that is both horizontal and $G$-invariant is called \emph{basic}.
When the $G$-action is understood,
we denote the set of basic $k$-forms on $M$ by $\Omega^k_{\rm basic}(M)$.
\end{Definition}

\begin{Remark} \label{r:wedge and d of basic}
The space of basic dif\/ferential forms on a $G$-manifold $M$
is closed under linear combinations, wedge products,
and exterior derivatives.
\end{Remark}

\begin{Remark}
Given a quotient map $\pi \colon X \to X'$ of dif\/feological spaces,
(more generally, given a so-called subduction,)
Iglesias-Zemmour~\cite[Section~6.38]{iglesias}
def\/ines a~``basic form'' to be a~dif\/ferential form on~$X$
that satisf\/ies the technical condition that appears in
Proposition~\ref{p:pullbackimgchar} above.
Our results show that, for smooth Lie group actions
where the identity component acts properly,
Iglesias-Zemmour's def\/inition agrees with the usual one.
\end{Remark}

Let $G$ be a Lie group, $H$ a closed subgroup, and $V$ a vector space
with a linear $H$-action. The equivariant vector bundle
\begin{gather*} G \times_H V \end{gather*}
over $G/H$ is obtained as the quotient of $G \times V$
by the anti-diagonal $H$-action $h \cdot (g,v) = (g h^{-1} , h \cdot v)$.
The $G$-action on $G\times_H V$ is $g\cdot[g',v]=[gg',v]$.

\begin{Theorem}[slice theorem]
\label{t:slice theorem}
Let $G$ be a Lie group acting properly on a manifold $M$. Fix $x\in M$.
Let $H$ be the stabiliser of $x$, and let $V = T_x M / T_x(G\cdot x)$
be the normal space to the orbit $G\cdot x$ at $x$,
equipped with the linear $H$-action that is induced
by the linear isotropy action of $H$ on $T_xM$.
Then there exist a $G$-invariant open neighbourhood $U$ of $x$
and a $G$-equivariant diffeomorphism $F\colon U\to G\times_H V$
that takes~$x$ to~$[1,0]$.
\end{Theorem}

\section{The pullback injects into basic forms}\label{sec:injects}

In this section we prove the easy part of Theorem~\ref{t:main}:
for a Lie group $G$ acting on a manifold $M$
with quotient map $\pi \colon M \to M/G$,
the pullback map $\pi^*$ is an injection
from the set of dif\/feological dif\/ferential forms on $M/G$
into the set of basic forms on $M$,
and $\pi^*$ is a dif\/feomorphism onto its image.

The following lemma is a special case of~\cite[Section~6.39]{iglesias}.

\begin{Lemma} \label{l:pullback is injection}
Let a Lie group $G$ act on a manifold $M$,
and let $\pi \colon M \to M/G$ be the quotient map.
Then the pullback map on forms,
$\pi^* \colon \Omega^k(M/G) \to \Omega^k(M)$, is an injection.
\end{Lemma}



\begin{proof}
More generally, let $X$ be a dif\/feological space, $\sim$ an equivalence
relation on $X$, and $\pi \colon X \to X/\!\!\sim$ the quotient map.
Then the pullback map on forms,
$\pi^* \colon \Omega^k(X/\!\!\sim) \to \Omega^k(X)$, is an injection.

By \hypref{r:linear}{Remark} it is enough to show that
the kernel of the pullback map
\begin{gather*}\pi^* \colon \ \Omega^k(X/\!\!\sim) \to \Omega^k(X)\end{gather*}
is trivial.
Let $\alpha \in \Omega^k(X/\!\!\sim)$
be such that $\pi^* \alpha = 0$.
Then, for any plot $p \colon U \to X$, we have $p^* \pi^* \alpha= 0$.
By the def\/inition of the quotient dif\/feology, this implies that
for any plot $q \colon U \to X/\!\!\sim$ we have $q^* \alpha = 0$.
Hence, $\alpha = 0$, as required.
\end{proof}

\begin{Proposition}[pullbacks from the orbit space are basic]
\label{p:pullbacksarebasic}
Let a Lie group $G$ act on a~mani\-fold~$M$.
Let $\alpha = \pi^*\beta$ for some $\beta \in \Omega^k(M/G)$.
Then $\alpha$ is basic.
\end{Proposition}

\begin{proof}
To show that $\alpha$ is $G$-invariant, note that for every $g \in G$, because $\pi\circ g=\pi$,
\begin{gather*}
g^* \alpha = g^* \pi^* \beta = \pi^* \beta = \alpha.
\end{gather*}

If $\alpha$ is a zero-form (that is, a smooth function)
then $\alpha$ is automatically horizontal and we are done.
Next, we assume that $\alpha$ is a dif\/ferential form of positive degree,
and we show that $\alpha$ is horizontal.
By Lemma~\ref{facts}, if $x\in M$ and $v\in T_x(G\cdot x)$, then there exists
$\xi\in\g$ such that
\begin{gather*}v=\frac{d}{dt}\Big|_{t=0}\exp(t\xi) \cdot x .\end{gather*}
Let $A_x\colon G\to M$ be the map sending $g$ to $g\cdot x$. Then, $v=(A_x)_*(\xi)$.

Thus, to show that $\alpha$ is horizontal, it is enough to show that $A_x^*\alpha=0$ for all $x\in M$. Indeed, it follows from the following commutative diagram that $A_x^*\alpha=A_x^*\pi^*\beta=0$,
\begin{gather*}
\begin{gathered}[b]
\xymatrix{
 G \ar[rr]^{A_x} \ar[d]_{g \mapsto \star} && M \ar[d]^{\pi} \\
 \{ \star \} \ar[rr]_{\star \mapsto [x]} && M/G
}\\[-17pt]
\null
\end{gathered}\tag*{\qed}
\end{gather*}
 \renewcommand{\qed}{}
\end{proof}

\begin{Remark}
\hypref{p:pullbacksarebasic}{Proposition} can also be deduced
from the following lemma:
a dif\/ferential form $\alpha$ on $M$ is basic
if and only if its pullbacks under the maps $G \times M \to M$
given by the projection $(g,x)\mapsto x$
and by the action $(g,x) \mapsto g\cdot x$ coincide.
For a proof of this lemma see, for example,
\cite[Lemma~3.3]{watts-groupoids}.
\end{Remark}

\begin{Proposition}[pullbacks via quotient maps]
\label{p:injectivity}
Let a Lie group $G$ act on a manifold $M$,
and let $\pi \colon M \to M/G$ be the quotient map.
Then the pullback map
\begin{gather*} \pi^* \colon \ \Omega^k(M/G) \to \pi^*\Omega^k(M/G) \end{gather*}
is a diffeomorphism,
where the target space is equipped with the subset diffeology
induced \linebreak from~$\Omega^k(M)$.
\end{Proposition}

\begin{proof}
More generally,
let $X$ be a dif\/feological space, let $\sim$ be an equivalence relation
on $X$, and let $\pi \colon X \to X/\!\!\sim$ be the quotient map.
We will show that the pullback map
\begin{gather*} \pi^* \colon \ \Omega^k(X/\!\!\sim) \to \pi^*\Omega^k(X/\!\!\sim) \end{gather*}
is a dif\/feomorphism,
where the target space is equipped with the subset dif\/feology
induced \linebreak from~$\Omega^k(X)$.

Clearly, $\pi^*$ is surjective to its image.
For the injectivity of~$\pi^*$, see \hypref{l:pullback is injection}{Lemma}.
As noted in Part~\eqref{f:pullback is sm} of \hypref{r:formfacts}{Remark},
$\pi^*$ is smooth.
We wish to show that the inverse map
\begin{gather*} (\pi^*)^{-1} \colon \ \pi^*\Omega^k(X/\!\!\sim) \to \Omega^k(X/\!\!\sim) \end{gather*}
is also smooth.

Fix a plot $p \colon U \to \Omega^k(X)$
with image in $\pi^*\Omega^k(X/\!\!\sim)$.
We would like to show
that $(\pi^*)^{-1} \circ p \colon U \to \Omega^k(X/\!\!\sim)$
is a~plot of $\Omega^k(X/\!\!\sim)$.
By \hypref{d:D on forms}{Def\/inition} of the dif\/feology
on spaces of dif\/ferential forms,
we need to show, given any plot $r \colon W \to X/\!\!\sim$
with $W \subset \RR^n$, that
\begin{gather*} (u,w)\mapsto r^* \big( \big( (\pi^*)^{-1} \circ p \big) (u) \big)\big|_w \end{gather*}
is a map to $\bigwedge^k \RR^n$ that is smooth in $(u,w) \in U \times W$.
Here, $(\pi^*)^{-1}$ is restricted to the image of~$\pi^*$,
on which it is well def\/ined because~$\pi^*$ is injective.

It is enough to show smoothness locally.
For any point $w\in W$ there exist an open neighbourhood $V\subseteq W$ of $w$
and a plot $q \colon V\to X$ such that $r|_V = \pi\circ q$.
For all $v\in V$, we have
\begin{gather*} r^* \big( \big((\pi^*)^{-1}\circ p\big)(u) \big)|_v = q^*(p(u))|_v,\end{gather*}
which is smooth in $(u,v) \in U \times V$
by the def\/inition of the standard functional dif\/feology on~$\Omega^k(X)$.
And so we are done.
\end{proof}

\section{Quotient in stages}\label{sec:q in stages}

In this section we give a technical result that we use in the next section.
All quotients, subsets, and products are assumed to be equipped
with the quotient, subset, and product dif\/feologies.

On the image of any Lie group homomorphism $H \to G$
there exists a unique manifold structure
such that the inclusion map of the image into $G$ is an immersion;
this follows from the second item of Lemma~\ref{facts}.
By \emph{Lie subgroup} of $G$ we refer to such an image.
Thus, Lie subgroups are injectively immersed subgroups
that are not necessarily closed.

Recall that a Lie group $G$ acts \emph{properly} on a manifold $N$
if the map $G\times N\to N\times N$ sending $(g,x)$ to $(x,g\cdot x)$
is proper. The action is said to have \emph{constant orbit-type}
if all stabilisers are conjugate.
If a Lie group acts properly and with a constant orbit type,
then the quotient is a~mani\-fold and the quotient map is a f\/ibre bundle.

\begin{Lemma}[quotient in stages] \label{l:q in stages}
Let a Lie group $G$ act on a manifold $N$.
Let $K$ be a Lie subgroup of~$G$
that is normal in $G$.
Also consider the induced action of $G$ on the quotient~$N/K$.
\begin{enumerate}\itemsep=0pt
\item[$(i)$]
There exists a unique map $ e \colon N/G \to (N/K) / G $
such that the following diagram commutes:
\begin{gather} \label{diagram q in stages}
\begin{split}
   \xymatrix{
 N \ar[rr]^{\pi_K} \ar[d]_{\pi_G} && N/K \ar[d]^{\pi_{G/K}} \\
 N/G \ar[rr]^{e} &&  (N/K)/G
}
\end{split}
\end{gather}

\item[$(ii)$]
The map $e$ is a diffeomorphism.

\item[$(iii)$]
The pullback map
\begin{gather*} \pi_K^* \colon \ \Omega^*(N/K) \to \Omega^*(N) \end{gather*}
restricts to a bijection from $\image \pi_{G/K}^*$ onto $\image \pi_G^*$.

\item[$(iv)$]
Suppose that $K$ acts on $N$ properly and with a constant orbit-type,
so that~$N/K$ is a~mani\-fold and $N \to N/K$ is a fibre bundle.
Then the pullback map $\pi_K^*$ also restricts to a bijection
from $\Omega^*_{\rm basic}(N/K)$ onto $\Omega^*_{\rm basic}(N)$.
Consequently, if one of the inclusions
$\image \pi_{G/K}^* \subset \Omega^*_{\rm basic}(N/K)$
and
$\image \pi_G^* \subset \Omega^*_{\rm basic}(N)$
of Proposition~{\rm \ref{p:pullbacksarebasic}} is an equality,
then so is the other.
\end{enumerate}
\end{Lemma}

\begin{Remark} \label{r:splits}
The $G$-action on $N/K$ factors through an action of $G/K$.
The quotients $(N/K)/G$ and $(N/K)/(G/K)$ coincide,
as they are quotients of $N/K$ by the same equivalence relation.
If $K$ is closed in $G$ (so that $G/K$ is a Lie group)
and $N/K$ is a manifold, then $G$-basic forms coincide with $(G/K)$-basic forms
on $N/K$.
\end{Remark}


\begin{proof}[Proof of Lemma~\ref{l:q in stages}]
Because $\pi_K$ is $G$-equivariant, such a map $e$ exists.
Because $\pi_G$ is onto, such a map $e$ is unique.
Because the preimage under $\pi_K$ of a $G$-orbit in $N/K$
is a single $G$-orbit in $N$, the map $e$ is one-to-one.
Because the maps $\pi_{G/K}$ and $\pi_K$ are onto,
the map $e$ is onto.

Thus, the map $e$ is a bijection. To show that it is a dif\/feomorphism,
it remains to show,
for every parametrisation $p \colon U \to N/G$,
that $p$ is a plot of $N/G$ if and only if $e \circ p$ is a plot of $(N/K)/G$.

Fix a parametrisation,
\begin{gather*} p \colon \ U \to N/G.\end{gather*}

Suppose that $p$ is a plot of $N/G$. Let $u \in U$.
Then there exist an open neighbourhood~$W$ of~$u$ in~$U$
and a plot $q \colon W \to N$ such that $p|_W = \pi_G \circ q$.
The composition $\pi_K \circ q$ is a~plot of~$N/K$, and
\begin{gather*} \pi_{G/K} \circ \pi_K \circ q = e \circ \pi_G \circ q = e \circ p|_W.\end{gather*}
Because $u \in U$ was arbitrary,
this shows that $e \circ p$ is a plot of $(N/K)/G$.

Conversely, suppose that $e \circ p$ is a plot of $(N/K)/G$.
Let $u \in U$. By applying the def\/inition of the quotient dif\/feology
at $\pi_{G/K}$ and then at $\pi_K$,
we obtain an open neighbourhood $W$ of $u$ in $U$
and a plot $r \colon W \to N$ such that
\begin{gather*} e \circ p|_W = \pi_{G/K} \circ \pi_K \circ r .\end{gather*}
By~(\ref{diagram q in stages}) and by the choice of $r$,
\begin{gather*}e \circ \pi_G \circ r = \pi_{G/K} \circ \pi_K \circ r = e \circ p|_W.\end{gather*}

Because $e$ is one-to-one, this implies that $\pi_G \circ r = p|_W$.
Because $u \in U$ was arbitrary, this shows that $p$ is a plot of $N/G$.
This completes the proof that $e$ is a dif\/feomorphism.

Because the diagram~\eqref{diagram q in stages} commutes, $\pi_K^*$
takes $\image \pi_{G/K}^*$ into $\image \pi_G^*$.
Because $e$ is a~dif\/feomorphism, we can consider its inverse.
From the commuting diagram
\begin{gather*} \xymatrix{
 N \ar[rr]^{\pi_K} \ar[d]_{\pi_G} && N/K \ar[d]^{\pi_{G/K}} \\
  N/G &&   (N/K) / G \ar[ll]_{e^{-1}}
} \end{gather*}
we see that $\pi_K^*$ takes $\image \pi_{G/K}^*$ \emph{onto} $\image \pi_G^*$.
Because $\pi_K^*$ is one-to-one (by \hypref{l:pullback is injection}{Lemma}),
we have a bijection
$\pi_K^* \colon \image \pi_{G/K}^* \to \image \pi_G^*$.

Now suppose that $K$ acts on $N$ properly and with a constant orbit-type,
so that $N/K$ is a~mani\-fold and $\pi_K$ is a f\/ibre bundle.
Then we know that $\pi_K^*$ is a bijection from the dif\/ferential forms
on $N/K$ to the $K$-basic dif\/ferential forms on $N$.

Because $\pi_K$ is $G$-equivariant,
$\pi_K^*$ takes $G$-invariant forms on $N/K$
to $G$-invariant forms on $N$
and $G$-horizontal forms on $N/K$ to $G$-horizontal forms on $N$.
So we have an injection
$\pi_K^* \colon \Omega^*_{\rm basic}(N/K) \to \Omega^*_{\rm basic}(N)$.

Let $\alpha$ be a $G$-basic form on $N$. In particular $\alpha$ is $K$-basic,
so there exists a form $\beta$ on $N/K$ such that $\alpha = \pi_K^* \beta$.
Because $\pi_K^*$ is one-to-one and $\alpha$ is $G$-invariant,
$\beta$ is $G$-invariant.
Because $\pi_K$ is $G$-equivariant and $\alpha$ is $G$-horizontal,
$\beta$ is $G$-horizontal.
This completes the proof that the map
$\pi_K^* \colon \Omega^*_{\rm basic}(N/K) \to \Omega^*_{\rm basic}(N)$
is a bijection.
\end{proof}

\section{The pullback surjects onto basics forms} \label{sec:surjects}

The main result of this section is Proposition~\ref{p:geomquot},
in which we give conditions on an action of a Lie group $G$ on a manifold $M$
under which \emph{every} basic form~$\alpha$ is the pullback
of some dif\/feological form on the quotient.
By \hypref{p:pullbackimgchar}{Proposition},
we need to show that $p_1^*\alpha = p_2^*\alpha$ for every two plots
$p_1 \colon U \to M$ and $p_2 \colon U \to M$
such that
for each $u\in U$ there is some $g \in G$ such that $p_2(u) = g \cdot p_1(u)$.
If $g$ can be chosen to be a smooth function of $u$,
then it is easy to conclude that $p_1^*\alpha = p_2^*\alpha$
if $\alpha$ is basic:

\begin{Lemma}\label{l:smooth case}
Let a Lie group $G$ act on a manifold $M$.
Let $p_1, p_2 \colon U \to M$ be plots.
Suppose that $p_2(u) = a(u) \cdot p_1(u)$
for some smooth function $a \colon U \to G$.
Then for every $\alpha \in \Omega^k_{\rm basic}(M)$
we have $p_1^*\alpha = p_2^*\alpha$.
\end{Lemma}

\begin{proof}
Pick a point $u \in U$ and a tangent vector $v \in T_uU$.
Let $\xi_1 = (p_1)_*v $ $\left( \in T_{p_1(u)}M \right)$
and $\xi_2 = (p_2)_*v$ $\left( \in T_{p_2(u)}M \right)$.
Let $g = a(u)$.
The directional derivative $D_v a|_u$ of $a(\cdot)$
in the direction of~$v$ has the form $\eta \cdot g$ $\left( \in T_gG \right)$
for some Lie algebra element~$\eta$ (i.e., it is the right translation
of~$\eta$ by~$g$).
We then have that
\begin{gather*} \xi_2 = g \cdot \xi_1 + {\eta_M}|_{p_2(u)}, \end{gather*}
where $g \cdot \xi_1$ is the image of~$\xi_1$
under the dif\/ferential (push-forward) map
$ g_* \colon T_{p_1(u)}M \to T_{p_2(u)}M $,
and where $\eta_M$ is the vector f\/ield on $M$
that corresponds to~$\eta$;
in particular~$\eta_M$ is everywhere tangent to the $G$ orbits.

Applying this to vectors $v^{(1)}, \ldots, v^{(k)} \in T_uU$,
we get that
\begin{align*}
 (p_2^* \alpha)|_u \big( v^{(1)} , \ldots , v^{(k)} \big)
 &= \alpha|_{ p_2(u) } \big( \xi_2^{(1)} , \ldots ,
 \xi_2^{(k)} \big)
 \quad \text{ where } \xi_2^{(j)} := (p_2)_* v^{(j)} \\
 &= \alpha|_{ p_2(u) }
 \big( g \cdot \xi_1^{(1)} + \eta^{(1)}_M , \ldots ,
 g \cdot \xi_1^{(k)} + \eta^{(k)}_M \big) \\
 & \qquad\qquad\qquad \text{ where } \xi_1^{(j)} := (p_1)_* v^{(j)}
 \text{ and }
 D_{v^{(j)}} a|_u = \eta^{(j)} \cdot g \\
 &= \alpha|_{ p_2(u) } \big( g \cdot \xi_1^{(1)} , \ldots ,
 g \cdot \xi_1^{(k)} \big)
 \quad \text{ because $\alpha$ is horizontal } \\
 &= \alpha|_{ p_1(u) } \big( \xi_1^{(1)} , \ldots ,
 \xi_1^{(k)} \big)
 \quad \text{ because $\alpha$ is invariant } \\
 &= (p_1^* \alpha)|_u \big( v^{(1)} , \ldots , v^{(k)} \big).\tag*{\qed}
\end{align*}
 \renewcommand{\qed}{}
\end{proof}

Unfortunately, in applying \hypref{p:pullbackimgchar}{Proposition},
it might be impossible to choose~$g$
to be a smooth function of~$u$:

\begin{Example}[$\ZZ_2\circlearrowright\RR$]
\label{x:plots differ nonsmoothly}
Let $M=\RR$, let $G = \{ 1, -1 \}$ with $(\pm1) \cdot x = \pm x$,
and let $\pi \colon M \to M/G$ be the quotient map.
Consider the two plots
$p_1 \colon \RR\to M$ and $p_2 \colon \RR\to M$
def\/ined as follows:
\begin{gather*} p_1(t) := \begin{cases} -e^{-1/t^2}  & \text{if} \ \ t<0, \\
 0  & \text{if} \ \ t=0, \\
 e^{-1/t^2}  & \text{if} \ \ t>0,
\end{cases}\end{gather*}
and
\begin{gather*} p_2(t):= \begin{cases} -e^{-1/t^2}  & \text{if} \ \ t\neq 0, \\
 0  & \text{if} \ \ t=0.
\end{cases}\end{gather*}
Then $\pi \circ p_1 = \pi \circ p_2$.
However, for $t<0$ we have $p_1(t)=1\cdot p_2(t)$,
whereas for $t>0$ we have $p_1(t)=-1\cdot p_2(t)$,
and so the two plots
do not dif\/fer by a continuous function to~$G$
on any neighbourhood of $t=0$.
\end{Example}

Our proofs use the following lemma.

\begin{Lemma}\label{l:baire}
Let $U \subseteq \RR^n$ be an open set.
Let $\{ C_i \}$ be a $($finite or$)$ countable collection
of relatively closed subsets of $U$ whose union is $U$.
Then the union of the interiors, $\bigcup_i \inter(C_i)$,
is open and dense in~$U$.
\end{Lemma}

\begin{proof}
This is a consequence of the Baire category theorem.
\end{proof}

%%% Countable group action case
%%% and reduction of compact Lie group to identity component

We now prove Proposition~\ref{p:geomquot} in the special case of a f\/inite group action. In this case, basic dif\/ferential forms are simply invariant dif\/ferential forms, as the tangent space to an orbit at any point is trivial.

\begin{Proposition}[case of a f\/inite group] \label{p:finite}
Let $G$ be a finite group, acting on a manifold~$M$.
Then every basic form on~$M$ is the pullback of a~$($diffeological$)$ differential form on~$M/G$.
\end{Proposition}

\begin{proof}
Fix a basic dif\/ferential $k$-form $\alpha$ on~$M$.
By \hypref{p:pullbackimgchar}{Proposition},
it is enough to show the following:
if $p_1 \colon U \to M$ and $p_2 \colon U \to M$ are plots
such that $\pi \circ p_1 = \pi \circ p_2$,
then $p_1^* \alpha = p_2^*\alpha$ on~$U$.
Fix two such plots $p_1 \colon U \to M$ and $p_2 \colon U \to M$.
For each $g \in G$ let
\begin{gather*} C_g := \{ u \in U\,|\, g \cdot p_1(u) = p_2(u) \}.\end{gather*}
By continuity, $C_g$ is closed for each~$g$.
By our assumption on~$p_1$ and~$p_2$,
\begin{gather*} U = \bigcup_{g\in G} C_g. \end{gather*}
By \hypref{l:baire}{Lemma},
the set $\bigcup_{g \in G} \inter(C_g)$
is open and dense in $U$. Thus, by continuity, it is enough to show
that $p_1^* \alpha = p_2^*\alpha$ on $\inter(C_g)$ for each $g \in G$. This, in turn, follows from the facts that, for each $g\in G$, we have $g\circ p_1=p_2$ on $\inter(C_g)$ and $g^*\alpha=\alpha$.
\end{proof}

Our next result, contained in \hypref{l:conntodisc}{Lemma}
and preceded by \hypref{l:C closed}{Lemma},
is a generalisation of the case of a f\/inite group action:
it shows that the property that interests us
holds for a Lie group action
if it holds for the action of the identity component of that
Lie group, assuming that the action of the identity component is proper.


\begin{Lemma}\label{l:C closed}
Let $G$ be a Lie group, and let $G_0$ be its identity component.
Assume that $G$ acts on a manifold $M$ such that the restricted action
of $G_0$ on $M$ is proper. Then, for any $\gamma\in G/G_0$,
and for any two plots $p_1 \colon U \to M$ and $p_2 \colon U \to M$, the set
\begin{gather*} C_\gamma := \{ u \in U\,|\, \exists\, g \in\gamma\text{ such that }
 g\cdot p_1(u)=p_2(u)\}\end{gather*}
is $($relatively$)$ closed in $U$.
\end{Lemma}

\begin{proof}
Because the $G_0$ action on $M$ is proper, the set
\begin{gather*}
\Delta := \{ (m,m') \in M \times M \, | \,
\exists\, g_0 \in G_0 \text{ such that } g_0 \cdot m = m' \},
\end{gather*}
being the image of the proper map $(g_0,m) \mapsto (m,g_0 \cdot m)$,
is closed in $M \times M$.

Fix $g' \in \gamma$. Because $G_0$ is normal in $G$,
we can express $\gamma$ as the left coset $G_0 g'$, and we have
\begin{gather*} C_\gamma = \{ u \in U \, | \, \exists\, g_0 \in G_0 \text{ such that }
g_0 g' \cdot p_1(u) = p_2(u) \} .\end{gather*}
We conclude by noting that $C_\gamma$ is the preimage of the closed set $\Delta$
under the continuous map $U \to M \times M$
given by $u \mapsto ( g' \cdot p_1(u) , p_2(u) )$.
\end{proof}

\begin{Lemma} \label{l:conntodisc}
Let $G$ be a Lie group. Let $G_0$ be the identity component of~$G$.
Fix an action of~$G$ on a~mani\-fold~$M$.
Suppose that the restricted $G_0$-action is proper,
and suppose that every $G_0$-basic differential form on~$M$
is the pullback of a diffeological form on $M/G_0$.
Then every $G$-basic differential form on~$M$
is the pullback of a diffeological form on~$M/G$.
\end{Lemma}

\begin{proof}
Fix a $G$-basic form $\alpha$ on $M$.
Let $\pi \colon M\to M/G$ be the quotient map,
and let $p_1 \colon U\to M$ and $p_2 \colon U\to M$ be plots
such that $\pi\circ p_1=\pi\circ p_2$.
Fix $\gamma\in G/G_0$ and
$g' \in \gamma$, and def\/ine $C_\gamma$
as in \hypref{l:C closed}{Lemma}. Def\/ine $\tilde{p}_1 \colon U \to M$
as the composition
$g' \circ p_1$. This is a plot of $M$,
and for any $u \in \inter(C_\gamma)$ we have
$p_2(u) = g_0 \cdot\tilde{p}_1(u)$ for some $g_0 \in G_0$.
Consider the restricted action of $G_0$ on $M$.
Let $\pi_0 \colon M \to M/G_0$ be the corresponding quotient map.
Then the restrictions $\tilde{p}_1|_{\inter(C_\gamma)}$
and $p_2|_{\inter(C_\gamma)}$ are plots of $M$,
and they satisfy
$\pi_0 \circ \tilde{p}_1|_{\inter(C_\gamma)}
 = \pi_0 \circ p_2|_{\inter(C_\gamma)}$.
By hypothesis, and because $\alpha$ is $G_0$-basic (as it is $G$-basic),
$\alpha$ is a pullback of a dif\/feological form on $M/G_0$.
By \hypref{p:pullbackimgchar}{Proposition},
$\tilde{p}_1^* \alpha = p_2^*\alpha$ on $\inter(C_\gamma)$.

But on $\inter(C_\gamma)$ we have
$\tilde{p}_1^* \alpha = p_1^* {g'}^* \alpha = p_1^* \alpha$
(since $\alpha$ is $G$-invariant),
and so $p_1^* \alpha = p_2^* \alpha$ on $\inter(C_\gamma)$.
Since $\gamma\in G/G_0$ is arbitrary,
and $\bigcup_{\gamma\in G/G_0}\inter(C_\gamma)$ is open and dense in $U$
by Lemmas~\ref{l:baire} and \ref{l:C closed},
%
%
from continuity we have that $p_1^*\alpha=p_2^*\alpha$ on all of $U$.
Finally, by \hypref{p:pullbackimgchar}{Proposition},
$\alpha$ is the pullback of a form on $M/G$.
\end{proof}


%%% Pointwise result at origin for connected compact linear orthogonal actions.

We proceed with two technical lemmas that we will use
to handle non-trivial compact connected stabilisers.

\begin{Lemma}\label{l:techn1}
Let $G$ be a compact connected Lie group acting orthogonally
on some Euclidean space $V = \RR^N$.
Let $g \in G$ and $\eta \in \g$ be such that $\exp(\eta) = g$.
Let $v \in V$. Then there exists $v' \in V$
such that $|v'| \leq |v|$ and $g \cdot v - v = \eta \cdot v'$.
\end{Lemma}

\begin{proof}
Since $V$ is a vector space, we identify tangent spaces at points of $V$
with $V$ itself.
We also identify elements of the group $G$ and of the Lie algebra $\g$
with the matrices by which these elements act on $V=\RR^N$
\begin{gather*}
 g \cdot v - v = \exp(t\eta) \cdot v \big|_0^1\\
\hphantom{g \cdot v - v}{}
 = \int_0^1 \left( \frac{d}{dt} \exp(t\eta) \cdot v\right) dt\\
\hphantom{g \cdot v - v}{}
 = \int_0^1 \left( \eta \cdot \exp(t\eta) \cdot v\right) dt \\
 \hphantom{g \cdot v - v}{}
 = \eta \cdot \int_0^1 \left( \exp(t\eta) \cdot v \right) dt.
\end{gather*}
So def\/ine $v':=\int_0^1\left(\exp(t\eta)\cdot v\right)dt$. Finally,
\begin{gather*}
 |v'| = \left| \int_0^1 \left( \exp(t\eta) \cdot v \right) dt \right|\\
\hphantom{|v'|}{}
 \leq \int_0^1 \left| \exp(t\eta) \cdot v\right| dt\\
 \hphantom{|v'|}{}
 = \int_0^1 |v| dt \quad \text{because the action is orthogonal} \\
\hphantom{|v'|}{} = |v|.
\end{gather*}
This completes the proof.
\end{proof}

\begin{Lemma}\label{l:techn2}
Let $G$ be a compact connected Lie group acting orthogonally
on some Euclidean space $V = \RR^N$.
Let $\gamma_1$ and $\gamma_2$ be smooth curves from $\RR$ into $V$
such that $ \gamma_1(0) = \gamma_2(0) = 0 $
and such that for every $t \in \RR$ there exists $g_t \in G$
satisfying $\gamma_2(t) = g_t \cdot \gamma_1(t)$.
Let $\xi_1 = \dot{\gamma}_1(0)$ and $\xi_2 = \dot{\gamma}_2(0)$.
Then, for every horizontal form $\alpha$ on $V$,
we have $(\xi_2-\xi_1)\hook\alpha|_0=0$.
\end{Lemma}

Note that, in the assumptions of this lemma,
$t \mapsto g_t$ is not necessarily continuous.

\begin{proof}
We claim that there exists a sequence of vectors $v'_n$
converging to $0$ in $V$, and a sequence~$\mu_n$ in~$\g$, such that
\begin{gather*} \xi_2 - \xi_1 = \underset{n\to\infty}{\lim} \mu_n\cdot v'_n. \end{gather*}

Indeed,
choose any sequence of non-zero real numbers $t_n$ converging to $0$.
For
each $n$ choose $\eta_n \in \g$ such that $\exp(\eta_n)=g_{t_n}$.
Since we are working on a vector space, we can subtract the curves
and consider $\gamma_2(t) - \gamma_1(t)$. We have
\begin{gather*}
 \xi_2 - \xi_1 = \frac{d}{dt} \Big|_{t=0} (\gamma_2(t)-\gamma_1(t))\\
  \hphantom{\xi_2 - \xi_1}{}
 = \underset{t \to 0} {\lim}
 \left( \frac{\gamma_2(t)-\gamma_1(t)}{t} \right)\\
 \hphantom{\xi_2 - \xi_1}{}
 = \underset{n\to\infty} {\lim}
 \left( \frac{\gamma_2(t_n)-\gamma_1(t_n)}{t_n} \right) \\
 \hphantom{\xi_2 - \xi_1}{}
 = \underset{n\to\infty}{\lim}
 \left( \frac{g_{t_n} \cdot \gamma_1(t_n) - \gamma_1(t_n)}{t_n} \right)\\
 \hphantom{\xi_2 - \xi_1}{}
 = \underset{n\to\infty}{\lim}
 \left( \frac{ \eta_n \cdot v'_n }{t_n} \right)
\end{gather*}
for some $\eta_n \in \g$ and
for some $v'_n \in V$ that satisfy $ |v'_n| \leq |\gamma_1(t_n)| $;
the last equality is a result of \hypref{l:techn1}{Lemma}.
Because $|v_n'| \leq |\gamma_1(t_n)| \underset{n\to\infty}{\to} |\gamma_1(0)|
 = 0$,
the claim holds with $\mu_n := \eta_n/t_n$.

We now have
\begin{gather*}
 (\xi_2-\xi_1) \hook \alpha|_0 = \underset{n\to\infty}{\lim}
 \left( (\mu_n\cdot v'_n) \hook(\alpha|_{v'_n}) \right)\\
 \hphantom{(\xi_2-\xi_1) \hook \alpha|_0}{}
 = \underset{n\to\infty}{\lim}
 \left( (\mu_n)_V\hook\alpha|_{v'_n} \right),
\end{gather*}
where $(\mu_n)_V$ is the vector f\/ield on $V$ induced by $\mu_n\in\g$.
Because $\alpha$ is horizontal, the last term above vanishes.
\end{proof}

The f\/inal ingredient that we need for Proposition~\ref{p:geomquot} is
the following property of the model that appears in the slice theorem
(Theorem~\ref{t:slice theorem}).
Let $G$ be a Lie group, $H$ a closed subgroup, and $V$ a vector space
with a linear $H$-action. Recall that the equivariant vector bundle
$ G \times_H V $
over $G/H$ is obtained as the quotient of $G \times V$
by the anti-diagonal $H$-action $h \cdot (g,v) = (g h^{-1} , h \cdot v)$
and that the $G$-action on $G\times_H V$ is $g\cdot[g',v]=[gg',v]$.

\begin{Lemma} \label{l:reduce to slice}
Suppose that every $H$-basic form on $V$ is the pullback
of a diffeological form on~$V/H$.
Then every $G$-basic form on $G \times_H V$ is the pullback
of a diffeological form on \mbox{$(G \times_H V)/G$}.
\end{Lemma}

\begin{proof}
Let $G \times H$ act on $G \times V$
where $G$ acts by left multiplication on the f\/irst factor
and where $H$ acts by the anti-diagonal action
$ h \colon (g,v) \mapsto (gh^{-1} , h \cdot v)$.
We have two maps:
\begin{gather*} %\label{d:triangle}
\xymatrix{
 & G \times V \ar[dl]_{\pr_2} \ar[dr]^{\pi_H} & \\
 V & &  \quad G \times_H V
}
\end{gather*}
Here, the map $\pi_H \colon G \times V \to G \times_H V$ is the quotient
by the $H$-action, and the projection to the second factor
$\pr_2 \colon G \times V \to V$
can be identif\/ied with the quotient by the $G$-action.

We apply quotient in stages (\hypref{l:q in stages}{Lemma})
in two ways:
taking the quotient
by~$G$ and then by~$H$, and taking the quotient by~$H$ and then by~$G$.
This gives the following commuting diagram:
\begin{gather} \label{two quotients}
\begin{split}
&
\xymatrix{
 & G \times V \ar[dl]_{\pr_2} \ar[dd]^{\pi_{G \times H}} \ar[dr]^{\pi_H} & \\
 V \ar[d]_{\pi_V} & & G \times_H V \ar[d]^{\pi} \\
   V/H \ar[r]^(0.3){e} & (G \times V) / (G \times H)
 &  (G \times_H V)/G \ar[l]_(0.42){e'}
}
\end{split}
\end{gather}
where $e$ and $e'$ are dif\/feomorphisms, and $\pi_V$, $\pi_{G\times H}$, and $\pi$ are the quotient maps.

By hypothesis, the inclusion
$\image \pi_V^* \subset \Omega^*_{\text{basic}}(V)$ is an equality.
Applying Part (iv) of \hypref{l:q in stages}{Lemma} to the left hand side
of the diagram~\eqref{two quotients},
we conclude that the inclusion
$\image \pi_{G \times H}^* \subset
 \Omega^*_{\text{basic}}(G \times V)$ is an equality.
Applying Part (iv) of \hypref{l:q in stages}{Lemma} to the right hand side
of the diagram~\eqref{two quotients},
we further conclude that the inclusion
$\image \pi^* \subset \Omega^*_{\text{basic}} (G \times_H V)$
is an equality.
\end{proof}

We are now ready to prove the main result of this section.

\begin{Proposition}[pullback surjects to basic forms]
\label{p:geomquot}
Let a Lie group $G$ act on a~mani\-fold~$M$.
Assume that the identity component of~$G$ acts properly.
Let $\pi \colon M \to M/G$ be the quotient map.
Then every basic form on $M$ is the pullback of a $($diffeological$)$ differential
form on $M/G$ via $\pi^*$.
\end{Proposition}

\begin{proof}
By \hypref{l:conntodisc}{Lemma}, to prove this result
for an arbitrary Lie group, it is enough to prove it
for the action of the identity component of the group.

For every non-negative integer $d$, consider the following two statements.
\begin{enumerate}\itemsep=0pt
\item[A(d):]
For every connected Lie group $K$ with $\dim K = d$,
and for every $K$-manifold $N$ on which the $K$-action is proper,
every $K$-basic form on~$N$ is the pullback of a dif\/ferential form on~$N/K$.

\item[B(d):]
For every Lie group $K$ with $\dim K = d$,
and for every $K$-manifold $N$ on which the identity component of $K$
acts properly, every $K$-basic form on~$N$ is the pullback
of a dif\/ferential form on $N/K$.
\end{enumerate}

Since the result holds for the trivial group, Statement A(0) is true.
By \hypref{l:conntodisc}{Lemma}, it follows that Statement B(0) is true.
Proceeding by induction, we f\/ix a positive integer $d$,
we assume that Statement~B(d$'$) is true for all $d' < d$,
and we would like to prove that Statement~B(d) is true.
By \hypref{l:conntodisc}{Lemma}, Statement~A(d) implies Statement~B(d);
thus, it is enough to prove that Statement~A(d) is true.
That is, we may now restrict to the special case of the proposition
in which the group is connected and the action is proper,
while assuming that the general case of the proposition is true
for all Lie groups of smaller dimension.

Now, let $G$ be a connected Lie group,
and let $M$ be a $G$-manifold on which the $G$-action is proper.
Fix a $G$-basic form~$\alpha$ on~$M$.
We would like to show that~$\alpha$ is the pullback of a~dif\/ferential form
on~$M/G$.

By \hypref{p:pullbackimgchar}{Proposition} we need to show,
for any two plots $p_1 \colon W \to M$
and $p_2 \colon W\to M$ for which $\pi \circ p_1 = \pi \circ p_2$,
that $p_1^* \alpha = p_2^* \alpha$.
Let $p_1$ and $p_2$ be two such plots.
Fix $u\in W$. We would like to show that $p_1^*\alpha|_u = p_2^*\alpha|_u$.

Let $x = p_2(u)$. Let $H$ be the stabiliser of $x$.
By Theorem~\ref{t:slice theorem} %the \hyperref[t:slice theorem]{slice theorem}
there exists
a $G$-invariant open neighbourhood $U$ of $x$ and an equivariant
dif\/feomorphism $F \colon U \to G \times_H V$
where $V = T_xM / T_x(G \cdot x)$.
Because $F$ is an equivariant dif\/feomorphism and $\alpha$ is $G$-basic,
$(F^{-1})^* \alpha$ is $G$-basic on $G \times_H V$.

Either $x$ is a f\/ixed point, or $x$ is not a f\/ixed point.

Suppose that $x$ is a f\/ixed point. Then $H=G$.
So $p_1(u) = p_2(u) = x$, and $F$ identif\/ies $U$ with $V = T_xM$,
sending $x$ to $0 \in V$.
Fixing a $G$-invariant Riemannian metric, we have
that $G$ acts linearly and orthogonally on $V$. Let $v \in T_uW$. Applying \hypref{l:techn2}{Lemma} to the curves
$\gamma_1(t) := F(p_1(u+tv))$ and $\gamma_2(t) := F(p_2(u+tv))$ in $V$
and to the basic form $(F^{-1})^*\alpha$ on $V$,
we obtain that $\dot{\gamma}_1(0) \hook (F^{-1})^*\alpha|_0
 = \dot{\gamma}_2(0) \hook (F^{-1})^*\alpha|_0$.
This, in turn, implies that $v \hook p_1^*\alpha|_u = v \hook p_2^* \alpha|_u$.
Because $v \in T_uW$ is arbitrary, we conclude
that $p_1^*\alpha|_u = p_2^*\alpha|_u$, as required.

Suppose that $x$ is not a f\/ixed point.
Then the stabiliser $H$ of $x$ is a proper subgroup of $G$.
Since $G$ is connected, $\dim H < \dim G$.
By the induction hypothesis, every $H$-basic form on $V$
is the pullback of a dif\/ferential form on $V/H$.
By \hypref{l:reduce to slice}{Lemma},
every $G$-basic form on $G \times_H V$
is the pullback of a dif\/ferential form on $(G \times_H V)/G$.
Because $F$ is an equivariant dif\/feomorphism,
every $G$-basic form on $U$
is the pullback of a dif\/ferential form on $U/G$.
So $\alpha|_U$ is the pullback of a dif\/ferential form on $U/G$.
This implies that $p_1^* \alpha|_u = p_2^* \alpha|_u$, as required.
\end{proof}

\begin{proof}[Proof of Theorem \ref{t:main}]
By \hypref{l:pullback is injection}{Lemma}
and \hypref{p:pullbacksarebasic}{Proposition},
the pullback is an injection into the space of basic forms.
By \hypref{alg iso}{Remark} and \hypref{p:injectivity}{Proposition},
as a map to its image,
the pullback is an isomorphism of dif\/ferential graded algebras
and a dif\/feological dif\/feomorphism.
By \hypref{p:geomquot}{Proposition}, if the identity component of $G$
acts properly, the image is the space of basic forms.
\end{proof}
\begin{thebibliography}{99}
\bibitem[1]{BH}
Baez J.C., Hof\/fnung A.E., Convenient categories of smooth spaces,
 \href{http://dx.doi.org/10.1090/S0002-9947-2011-05107-X}{\textit{Trans. Amer. Math. Soc.}} \textbf{363} (2011), 5789--5825,
 \href{http://arxiv.org/abs/0807.1704}{arXiv:0807.1704}.

\bibitem[9]{DK}
Duistermaat J.J., Kolk J.A.C., Lie groups, \href{http://dx.doi.org/10.1007/978-3-642-56936-4}{\textit{Universitext}}, Springer-Verlag,
 Berlin, 2000.

\bibitem[10]{GGK}
Guillemin V., Ginzburg V., Karshon Y., Moment maps, cobordisms, and
 {H}amiltonian group actions, \href{http://dx.doi.org/10.1090/surv/098}{\textit{Mathematical Surveys and Monographs}},
 Vol.~98, Amer. Math. Soc., Providence, RI, 2002.

\bibitem[14]{iglesias}
Iglesias-Zemmour P., Dif\/feology, \href{http://dx.doi.org/10.1090/surv/185}{\textit{Mathematical Surveys and Monographs}},
 Vol.~185, Amer. Math. Soc., Providence, RI, 2013.

\bibitem[15]{IK}
Iglesias-Zemmour P., Karshon Y., Smooth {L}ie group actions are parametrized
 dif\/feological subgroups, \href{http://dx.doi.org/10.1090/S0002-9939-2011-11301-7}{\textit{Proc. Amer. Math. Soc.}} \textbf{140} (2012),
 731--739, \href{http://arxiv.org/abs/1012.0107}{arXiv:1012.0107}.

\bibitem[17]{koszul}
Koszul J.L., Sur certains groupes de transformations de {L}ie, in G\'eom\'etrie
 dif\/f\'erentielle ({C}olloques {I}nternationaux du {C}entre {N}ational de la
 {R}echerche {S}cientif\/ique, {S}trasbourg, 1953), Centre National de la
 Recherche Scientif\/ique, Paris, 1953, 137--141.

\bibitem[19]{palais}
Palais R.S., On the existence of slices for actions of non-compact {L}ie
 groups, \href{http://dx.doi.org/10.2307/1970335}{\textit{Ann. of Math.}} \textbf{73} (1961), 295--323.

\bibitem[29]{watts-groupoids}
Watts J., The orbit space and basic forms of a proper Lie groupoid,
 \href{http://arxiv.org/abs/1309.3001}{arXiv:1309.3001}.

\end{thebibliography}
\end{document}
